\documentclass[11pt]{article}

\usepackage{amssymb,amsmath,amsfonts,amsthm}
\usepackage{latexsym}
\usepackage{graphics}
\usepackage{indentfirst}
\usepackage{lmodern}
\usepackage{hyperref}
\usepackage{mathrsfs}
\usepackage{array}
\usepackage{mathtools}
\usepackage{graphicx}
\usepackage{setspace}
\usepackage{enumerate}
\usepackage{tikz}
\usepackage{url,color}
\usepackage{chngcntr}
\usepackage{multicol}
\usepackage{comment}
\usepackage{float}
\usepackage{algpseudocode}
\usepackage{subcaption}
\usepackage[ruled]{algorithm2e}
\usepackage{booktabs}
\usepackage[numbers,sort&compress]{natbib}

\AtBeginEnvironment{algorithm}{\setstretch{1.2}}

\setlength{\textwidth}{15.5cm} \setlength{\headheight}{0.5cm} \setlength{\textheight}{21.5cm}
\setlength{\oddsidemargin}{0.25cm} \setlength{\evensidemargin}{0.25cm} \setlength{\topskip}{0.5cm}
\setlength{\footskip}{1.5cm} \setlength{\headsep}{0cm} \setlength{\topmargin}{0.5cm}


\newtheorem{theorem}{Theorem}[section]
\newtheorem{lemma}{Lemma}[section]
\newtheorem{proposition}{Proposition}[section]
\newtheorem{corollary}{Corollary}[section]
\newtheorem{definition}{Definition}[section]
\newtheorem{remark}{Remark}[section]

\newcommand{\N}{\mathbb{N}}
\newcommand{\Z}{\mathbb{Z}}
\newcommand{\E}{\mathbb{E}}
\newcommand{\R}{\mathbb{R}}
\newcommand{\Var}{\mathrm{Var}}

\allowdisplaybreaks


\begin{document}


%%%%%%%%%%%%%%%%%%%%%%%%%    Title and Authors
\title{Optimization-Based Decision Support for Equitable Expansion of Federally Qualified Health Centers}

\author{Kim Erwin$^2$, Hemanshu Kaul$^1$, Alaittin K\i rt\i \c{s}o\u{g}lu$^1$}

\date{}

\footnotetext[1]{Department of Applied Mathematics, Illinois Institute of Technology, Chicago, IL 60616. E-mail: {\tt {kaul@iit.edu, kirtisoglu@hawk.iit.edu}}}
\footnotetext[2]{Institute of Design, Illinois Institute of Technology, Chicago, IL 60616. E-mail: {\tt {kerwin@id.iit.edu}}}
\maketitle

%%%%%%%%%%%%%%%%%%%%%%%%%    Abstract
\begin{abstract}
TODO
\end{abstract}
\noindent{\bf Keywords:} Primary care access, Medically Uniderserved Areas, Health Equity, Federally Qualified Health Centers

%%%%%%%%%%%%%%%%%%%%%%%%%    Contents
\tableofcontents

\newpage



%%%%%%%%%%%%%%%%%%%%%%%%%%%%%%%%%%%%%%%%%%%%%%%%%%
\section{Introduction}
%%%%%%%%%%%%%%%%%%%%%%%%%%%%%%%%%%%%%%%%%%%%%%%%%%

% 1. Why primary care and community health centers matter
The community health center program presents an opportunity to address critical public health 
vulnerabilities resulting from a lack of primary care services. The health center model has 
proven to increase access to primary care, reduce health disparities, lower costs, and improve 
community health. The number of clinic sites operated by these organizations has increased from 
less than 10,000 in 2015 to over 14,000 in 2021, serving 6 million additional patients (a 24\% increase) 
over the same period \textcolor{red}{[CITE: NACHC or HRSA annual report for clinic growth statistics]}.

In addition to providing direct services, health centers serve as the training ground for the next
generation of primary care, behavioral health, and oral health providers \textcolor{red}{[CITE: NACHC or HRSA workforce report]}. A lack of access to primary
care can weaken the public health response to future pandemics and could have dire consequences for
patients, especially those facing compounding access barriers \cite{NASEM2021}. Patients with limited access to
transportation, financial limitations, and other access barriers may forgo care rather than travel
long distances to see a primary care provider.

% 2. Medically underserved areas and the designation gap
Medically underserved refers to people who live in a federally designated Health Professional 
Shortage Area (HPSA) or Medically Underserved Area (MUA) or belong to a Medically Underserved 
Population (MUP). Since 2014, the number of medically disenfranchised people has nearly doubled 
from 56 million to over 100 million. A study by HealthLandscape at the American Academy of Family 
Physicians, commissioned by NACHC, found that over 100 million people in the United States may not 
have access to a usual source of primary care due to a shortage of providers in their community \cite{NACHC2014}.

The Health Resources and Services Administration (HRSA) designates MUAs based on the Index of Medical 
Underservice (IMU), which incorporates four indicators: physician supply, poverty rate, infant mortality, 
and the proportion of elderly residents \cite{HRSA_BPHC}. Areas with an IMU score of 62 or less qualify 
for MUA designation and become eligible for the establishment of Federally Qualified Health Centers (FQHCs). 
Nationally, \$3.7 billion in federal dollars were allocated for FQHCs in 2014 \cite{NACHC2014}.

However, prior research has identified an important inconsistency: some geographic areas have IMU scores 
below the eligibility threshold but are not formally designated as MUAs \cite{kim2020uneven}. This discrepancy
 creates a gap between measured healthcare need and formal designation, potentially affecting where federal 
 resources are directed. MUA designation usually increases with poverty, but it actually declines once a 
 neighborhood hits a certain tipping point of extreme poverty and racial segregation. Highly underserved 
 areas often lack the community organization and resources required to successfully navigate the federal
 application process for MUA status \cite{kim2020uneven}.

% 3. Chicago context
Some very alarming statistics reveal population-level disparities within our social and health-care systems. 
One prominent example is the 30-year gap in life expectancy between people living in the poor, predominantly 
African American neighborhoods on Chicago's South Side, as compared to those in the more affluent, predominantly 
white neighborhoods just nine miles away in Chicago's Loop---the largest life expectancy gap in the United States. 
Real solutions are only possible with equitable strategies that recognize the contribution of historical and 
contemporary policies, economics, and health-care access \cite{smithcommunity}.

In Chicago, 68 new clinics opened and only 12 closed since 2009, yet these additions often failed to improve 
access in the areas that needed them most. Most new clinics were not located in the most severely underserved 
South and West Side neighborhoods, leading to a safety net that perpetuates racial and socioeconomic inequality 
\cite{kim2020uneven}. Recent closures further strain the system: a Humboldt Park health clinic that served 
neighbors and low-income people for nearly a decade was closed \cite{Bloom2023}, Howard Brown Health closed 
two clinics amid financial pressure \cite{Schencker2024}, and Advocate Health Care closed nearly 50 clinics 
inside Walgreens stores \cite{Japsen2025}.

% 4. Gap in methodology
The process for opening new FQHC facilities typically occurs through grant applications from local organizations 
rather than through systematic spatial planning. HRSA does not use formal optimization models like location-allocation
or flow models \textcolor{red}{[CITE: source for HRSA process description]}. Community organizations apply for federal funding, and proposals are evaluated based on population
need, access gaps, financial viability, and organizational capacity. As a result, facility placement may not fully
reflect national or regional access optimization.

There is currently no large-scale optimization-based location-allocation model that integrates underserved area 
indicators, facility capacity constraints, population demand, and equity in geographic access. Most existing studies 
analyze healthcare access descriptively rather than optimizing facility placement. To the best of our knowledge, there is no single paper for Chicago providing actionable optimization-based strategies for decision makers.

% 5. Our contribution
We propose a flow-based spatial optimization framework for equitable planning of primary care access in Chicago. The framework:
\begin{itemize}
    \item Computes IMU scores at finer geographic projections (census blocks) to identify the MUA designation gap at the boundaries,
    \item Proposes alternative IMU formulations such as probabilistic scores and smooth boundary transitions to address the sharp cutoff at IMU $= 62$,
    \item Designs a hierarchical location-allocation model with a referral system linking FQHCs to hospitals,
    \item Uses CTA public transit travel times reflecting real access barriers,
    \item Uses FalCom, an MCMC sampler for hierarchical capacitated facility location, as the computational engine,
    \item Provides an interactive decision-support dashboard for community organizations
    and investors preparing FQHC grant applications: given a proposed facility location,
    the dashboard runs a FalCom chain and returns a predicted HRSA application score,
    IMU score with smooth boundary version, patient flow distribution, referral network
    structure with HRSA compliance check, and ensemble robustness statistics across the
    distribution of feasible partitions.
\end{itemize}


%%%%%%%%%%%%%%%%%%%%%%%%%%%%%%%%%%%%%%%%%%%%%%%%%%
\section{Background}\label{sec:background}
%%%%%%%%%%%%%%%%%%%%%%%%%%%%%%%%%%%%%%%%%%%%%%%%%%


\subsection{Federally Qualified Health Centers}

Health centers need additional resources to sustain program growth and expand access to primary 
care in underserved communities. New FQHCs are usually not built through a centralized optimization 
plan. Instead, the process works roughly as follows: (i) an area is identified as underserved using 
metrics such as the IMU or HPSA scores; (ii) community organizations apply for federal funding to 
open or expand a health center; (iii) HRSA evaluates proposals based on population need, access gaps, 
financial viability, and organizational capacity; and (iv) HRSA awards Section 330 grants to selected 
applicants \cite{HRSA_BPHC}.

FQHCs primarily provide primary care, preventive care, and chronic disease management, but they do not 
provide most specialized or inpatient services. Patients must be referred to hospitals, specialists, 
and diagnostic facilities. HRSA requires formal referral relationships, and research shows that strong 
referral relationships are critical for community health center effectiveness \textcolor{red}{[CITE: source for FQHC referral effectiveness]}.


\subsection{MUA Designation and the IMU Score}\label{sec:imu}

An MUA is a geographical area designated as lacking access to primary care services and eligible for enhanced 
federal support. The geographical area can be a whole county, a group of neighboring counties, or a group of 
neighboring census tracts. HRSA is responsible for the evaluation of eligibility and designation of MUAs based 
on the IMU score \cite{HRSA_BPHC}.

The IMU is computed as a weighted sum of four indicators, each with a standardized contribution:
\begin{enumerate}
    \item The ratio of primary care physicians per 1,000 residents,
    \item The infant mortality rate,
    \item The percentage of the population with incomes below the poverty level,
    \item The percentage of the population aged 65 years or older.
\end{enumerate}

\begin{align}
\text{IMU} &= S_{\text{phys}} + S_{\text{poverty}} + S_{65+} + S_{\text{IMR}} \label{eq:imu} \\[6pt]
S_{\text{phys}} &= 28.7 \cdot \frac{P_{\max} - P}{P_{\max} - P_{\min}} \in [0,\,28.7] \\[4pt]
S_{\text{poverty}} &= 25.1 \cdot \frac{V_{\max} - V}{V_{\max} - V_{\min}} \in [0,\,25.1] \\[4pt]
S_{65+} &= 20.2 \cdot \frac{A_{\max} - A}{A_{\max} - A_{\min}} \in [0,\,20.2] \\[4pt]
S_{\text{IMR}} &= 26.0 \cdot \frac{M_{\max} - M}{M_{\max} - M_{\min}} \in [0,\,26.0]
\end{align}

The total IMU score ranges from 0 to 100, where lower values indicate greater underservice. 
Areas with IMU $\leq 62$ qualify for MUA designation.

For metropolitan areas, the rational service area is defined as a group of census tracts that 
represent a neighborhood of homogeneous socio-economic and demographic characteristics. Although
DHHS shortage area designation methods take into account some non-spatial factors such as age
and socio-economic status, they are primarily regional availability measures that quantify the
distribution of supply versus demand within a predefined region, often expressed as a ratio of
population to practitioner \cite{ricketts2007designating}.

The advantage of such a regional availability approach is that it is simple, straightforward
to implement (the data for physicians and population are readily available and boundaries can
be easily located), and convenient to administer because the government infrastructure is already
in place \cite{ricketts2007designating}.


\subsection{Limitations of the IMU and the Designation Gap}\label{sec:imu-limitations}

There is significant literature discussing the limitations of IMU and MUA designation. The IMU formula
relies on weights created in the 1970s. Ricketts et al.\ \cite{ricketts2007designating} argue that the 
method no longer adequately captures modern healthcare access barriers. Healthcare utilization patterns 
and physician workforce distribution have changed substantially. Major determinants of healthcare access
such as insurance coverage, transportation barriers, racial disparities, chronic disease prevalence are 
not included in the IMU formula.

The designation rule IMU $\leq 62$ creates a sharp cutoff: two areas with IMU scores of 61 and 63 receive 
entirely different policy treatment, yet their healthcare needs may be nearly identical. Kim et al.\ \cite{kim2020uneven} 
show that some Chicago tracts meet IMU eligibility criteria but are not designated MUAs because of
administrative boundaries. Census tracts that are eligible for MUA status but not designated were predominantly 
Black and poor, located on the South Side of Chicago.

We address these limitations in two ways. First, we compute IMU scores at the census tract level across Chicago 
to identify tracts that are eligible but undesignated. Second, we propose alternative scoring approaches such as
probabilistic IMU scores and smooth boundary transition functions that replace the binary cutoff with a 
continuous measure of underservice (Section~\ref{sec:new-imu}).


\subsection{The Chicago Context}\label{sec:chicago}

Kim et al.\ \cite{kim2020uneven} used the spatial distribution of MUA designations and FQHCs in Chicago to evaluate 
differences between MUA-eligible neighborhoods that received designation and those that did not. Of Chicago's census 
tracts, 391 were in affluent neighborhoods not eligible for MUA designation, mostly on the north, northwest side, and 
downtown. The 390 tracts that were eligible and designated for MUA status were clustered on the far north and west side. 
The 87 tracts that were eligible (IMU scores 48.6--61.7) but not designated as MUAs were predominantly on the South Side.

\begin{table}[h]
\centering
\caption{Distribution of FQHCs in Chicago by MUA status \cite{kim2020uneven}.}
\label{tab:mua-distribution}
\begin{tabular}{lrrr}
\toprule
 & Chicago & MUA & Non-MUA \\
\midrule
Newly added clinics (since 2009) & 68 & 41 & 3 \\
FQHCs per 100,000 poor & 39.6 & 40.8 & 15.5 \\
Total number of FQHCs & 187 & 112 & 10 \\
Total population & 2,683,422 & 1,252,783 & 259,226 \\
\bottomrule
\end{tabular}
\end{table}

Neighborhoods with a greater number of community organizations per 1,000 residents were more likely to receive MUA 
designation, highlighting that administrative capacity, not just medical need drives designation outcomes \cite{kim2020uneven}.


\subsection{Travel Time and Transit Equity}\label{sec:travel}

In the city of Chicago, 36.9\% of African American households had no vehicle compared to 23.2\% of White households \textcolor{red}{[CITE: ACS or census vehicle ownership data]}.
Physical distance is therefore less important than public transit accessibility, as many underserved residents lack
private vehicles.

Ermagun and Tilahun \cite{ermagun2020equity} study the equity of access to different destinations by public transit 
in Chicago. They find that the far south, southwest, and northwest of the city are among areas with lowest accessibility. 
Looking at accessibility to hospitals, the largest drops associated with population increases are among the elderly, 
low-income, and families with children aged 5--19. Despite the need of these cohorts to have access to healthcare, 
public transit does not serve them well, while providing better access to areas with higher percentages of highly-educated 
people. The South Side, which is predominantly African American and low-income, is among the least accessible areas. 
Chicago is ranked 3rd nationally for greatest accessibility by transit and 4th for greatest accessibility by walking \textcolor{red}{[CITE: source for Chicago transit ranking]}.

In this work, we use CTA (Chicago Transit Authority) public transit travel times as the distance metric for our model, 
reflecting the actual transportation options available to the underserved population.


\subsection{HRSA Compliance Requirements}\label{sec:hrsa-compliance}

HRSA requires that health center service sites be accessible to the patient population relative to where the population 
lives or works \textcolor{red}{[CITE: HRSA Program Compliance Manual]}. Health centers must consider distance and travel 
time between service sites when assessing impact on patient access. HRSA also requires that health centers establish and 
maintain collaborative relationships with other health care providers within the catchment area, local hospitals, and 
specialty providers to:
\begin{itemize}
    \item Reduce non-urgent use of hospital emergency departments,
    \item Ensure continuity of care across community providers,
    \item Provide access to services not available through the health center.
\end{itemize}

Facility placement should therefore consider proximity to hospitals and referral networks. This motivates the hierarchical 
structure of our model, where each FQHC is assigned to a referral hospital.



%%%%%%%%%%%%%%%%%%%%%%%%%%%%%%%%%%%%%%%%%%%%%%%%%%
\section{Model and Formulation}\label{sec:model}
%%%%%%%%%%%%%%%%%%%%%%%%%%%%%%%%%%%%%%%%%%%%%%%%%%

\subsection{Doubly-Constrained Spatial Flow Model}\label{sec:flow-model}

Let $I$ denote the set of demand nodes (census tracts or blocks) and $J$ the set of open FQHCs
in a given spatial partition. Each demand node $i \in I$ carries demand $d_i$, representing
realized patient counts from the UDS or census-based population. Each FQHC $j \in J$ has
capacity $c_j$ derived from UDS workforce and visit data. The travel cost $t_{ij}$ is the
CTA public transit travel time from $i$ to $j$.

\paragraph{Distance decay.}
We model the friction of travel using an exponential decay function:
\begin{equation}
    f(t_{ij}) = e^{-\lambda t_{ij}}
    \label{eq:decay}
\end{equation}
where $\lambda > 0$ is a decay parameter calibrated to observed patient flows via the UDS data.

\paragraph{Flow model.}
Patient flows $x_{ij}$ from demand node $i$ to FQHC $j$ are given by:
\begin{equation}
    x_{ij} = A_i \cdot B_j \cdot d_i \cdot c_j \cdot f(t_{ij})
    \label{eq:flow}
\end{equation}
where $A_i$ and $B_j$ are balancing factors ensuring that both supply and demand constraints
are simultaneously satisfied.

\paragraph{Balancing factors.}
The factors $A_i$ and $B_j$ are the unique solution to:
\begin{align}
    A_i &= \frac{1}{\displaystyle\sum_{j \in J} B_j \cdot c_j \cdot f(t_{ij})}
    \label{eq:Ai} \\[6pt]
    B_j &= \frac{1}{\displaystyle\sum_{i \in I} A_i \cdot d_i \cdot f(t_{ij})}
    \label{eq:Bj}
\end{align}
These are solved via Furness (IPF) iteration: initialize $B_j = 1$ for all $j$, then alternate
updating $A_i$ via~\eqref{eq:Ai} and $B_j$ via~\eqref{eq:Bj} until convergence, typically
within 20--50 iterations. The resulting flows satisfy:
\begin{align}
    \sum_{j \in J} x_{ij} &= d_i \quad \forall\, i \in I
    \label{eq:demand-conservation} \\
    \sum_{i \in I} x_{ij} &= c_j \quad \forall\, j \in J
    \label{eq:capacity-conservation}
\end{align}
Constraint~\eqref{eq:demand-conservation} ensures all demand is allocated;
constraint~\eqref{eq:capacity-conservation} ensures each FQHC operates at full capacity.
A singly-constrained variant (demand side only, $\sum_i x_{ij} \leq c_j$) is also considered
for scenarios where FQHCs may operate below capacity.

\paragraph{Role of FalCom.}
FalCom generates the spatial partition — determining which census tracts fall within each
FQHC's catchment area and which FQHCs are assigned to which referral hospitals. Given a
partition, the flow model~\eqref{eq:flow}--\eqref{eq:Bj} is solved to compute realized flows
$x_{ij}$. The energy function evaluated at each MCMC step is a flow-based equity or access
score. A natural choice is the population-weighted average travel time:
\begin{equation}
    E = \frac{\displaystyle\sum_{i \in I}\sum_{j \in J} x_{ij} \cdot t_{ij}}{\displaystyle\sum_{i \in I} d_i}
    \label{eq:energy}
\end{equation}
Alternative energy functions include the Gini coefficient of tract-level access
$\bigl(\sum_j x_{ij} t_{ij}\bigr)/d_i$ or the fraction of demand within a target travel
time threshold. The FalCom sampler generates a distribution over partitions weighted by
$\mu_\beta \propto \exp(-\beta E)$, concentrating on equitable, low-energy configurations
while preserving sampling diversity.

\paragraph{Calibration.}
Because UDS provides both realized patient flows and facility capacity, $\lambda$
in~\eqref{eq:decay} can be estimated from observed data rather than chosen ad hoc. This is an
empirical advantage over most facility location studies, which rely on synthetic demand or
proxy travel costs.


\subsection{Hierarchical Catchment Areas}

In the context of Chicago, a hierarchical system divides facilities into levels: primary care (Level~1) and hospital/specialty 
care (Level~2). Tao et al.\ \cite{tao2020hierarchical} propose a hierarchical 2SFCA (H2SFCA) method for measuring spatial 
accessibility to hierarchical facilities, considering varied catchment area sizes, distance decay effects, and transport modes 
at each level. Wan et al.\ \cite{wan2012three} develop a three-step floating catchment area method that indicates strong potential 
for identifying health professional shortage areas.


\subsection{Redesigning the IMU Score}\label{sec:new-imu}

The IMU formula is not merely outdated; it is structurally flawed in ways that
systematically misallocate resources. In this section we identify the specific
failures, review prior reform attempts and the political constraints that blocked
them, and propose a redesigned scoring framework that integrates with the
optimization model.


\subsubsection{Why the Formula Is Wrong}\label{sec:imu-wrong}

Beyond the well-documented omission of insurance, transportation, and chronic
disease burden \cite{ricketts2007designating}, the IMU has deeper structural
problems.

\paragraph{Miscalibrated conversion tables.}
Each indicator is mapped to a weighted score via lookup tables calibrated to
1970s national distributions. The national poverty rate, provider supply,
infant mortality rate, and age distribution have all shifted substantially.
A census tract that would have scored $10/25.1$ on poverty in 1975 may score
$18/25.1$ today on the same table, not because it worsened but because the
scale no longer matches the distribution. The scores do not measure what they
claim to measure.

\paragraph{Statistical noise in infant mortality.}
A typical census tract has 300--500 births per year. At this sample size, a
single infant death changes the rate by 2--3 per thousand, which can shift
the IMU by several points. The indicator carries enormous variance and almost
no signal at the geographic level where IMU is computed. Funding decisions are
made on noise.

\paragraph{Linear independence assumption.}
The four indicators are summed with fixed weights, implying that they are
independent and equally substitutable. A tract with no providers but low poverty
scores the same as a tract with adequate providers but extreme poverty. In
reality these factors interact: poverty without providers is catastrophically
worse than either alone. A linear sum cannot capture compounding disadvantage.

\paragraph{Scale compression at the extremes.}
The $0$--$25$ mapping for each indicator flattens differences among the most
underserved areas. A tract with zero providers and a tract with 5 per 100,000
may both map to 24 or 25. For optimization, this is fatal: the score must
differentiate precisely where need is greatest. The bounded sum (0 to 100)
saturates at both tails, preventing the score from reflecting how far below
the threshold an area truly falls.

\paragraph{Provider counts ignore capacity.}
The physician-to-population ratio counts heads, not capacity. A tract with one
physician seeing 500 patients and a tract with one physician seeing 3,000
patients receive the same ratio. Our provider matching reveals that 34 of 166
FQHC sites in Chicago have zero providers registered at their address in the
NPPES, yet the IMU formula would still count those sites as serving their area.


\subsubsection{Prior Reform Attempts}\label{sec:imu-reform}

Ricketts et al.\ \cite{ricketts2007designating} proposed an alternative
methodology, the Index of Primary Care Underservice, developed at the Cecil G.\
Sheps Center at UNC Chapel Hill. Their six-step method computes an age-sex
adjusted effective population, counts clinician FTEs (including nurse
practitioners and physician assistants, not only physicians), and applies a
barrier adjustment derived from regression on nine indicators: poverty below
200\% FPL, unemployment, racial minority share, Hispanic share, linguistic
isolation, infant mortality or low birthweight, standardized mortality ratio,
low population density, and share aged 65+. Each indicator is converted to
a national percentile rank and then to a partial score via regression-derived
weights. The final output is an adjusted population-to-clinician ratio.

HRSA adopted this methodology in a proposed rule published in the Federal
Register on February 29, 2008 (73 FR 11232--11281)
\cite{HRSA2008proposed}. The proposal would have unified the separate MUA
and HPSA designation systems into a single three-tier framework. HRSA
received overwhelmingly negative public comments and withdrew the rule
without finalizing it. Many areas that held MUA or HPSA designations would
have lost them under the new formula, threatening Section 330 funding for
existing FQHCs. A previous proposed rule in 1998 was similarly withdrawn.

Section 5602 of the Affordable Care Act (2010) required HRSA to establish a
new methodology via negotiated rulemaking. A 28-member committee voted 21--2
in favor of a revised approach in October 2011, but the Negotiated Rulemaking
Act required unanimity for HHS to publish the recommendation as an interim
rule. HHS did not act. The 1975 formula remains in effect today
\cite{HRSA2008proposed}.


\subsubsection{The Political Constraint}\label{sec:imu-political}

Every reform attempt has failed for the same reason: the IMU score is fused
to a binary funding gate. The chain is:
\[
\text{IMU} \leq 62 \;\longrightarrow\; \text{MUA designation}
\;\longrightarrow\; \text{Section 330 funding eligibility.}
\]
A one-point difference between 62 and 63 means the difference between
millions of dollars and zero. Any formula change that reshuffles scores
threatens existing clinics serving real patients. The binary cutoff makes the
formula politically unreformable: improving the measurement changes the money.


\subsubsection{The Dilution Problem}\label{sec:dilution}

If the city is partitioned into contiguous districts (as required for honest
catchment area design), each district's IMU is the population-weighted average
of its census tracts. Contiguity forces each district to include neighboring
tracts that may differ in need, pulling the average toward the middle. A
district straddling a high-poverty tract and a moderate-income tract may
average to an IMU of 63 or 64, just above the cutoff, even though half its
population is severely underserved.

This creates a tension between honest boundaries and the binary threshold:
\begin{itemize}
  \item Gerrymandering (the current system) allows applicants to draw
    boundaries that cherry-pick the worst tracts and push IMU below 62.
  \item Contiguous partitioning produces honest boundaries but dilutes
    concentrated pockets of need, causing more communities to fail the
    threshold.
\end{itemize}
The more honest the boundary-drawing process, the fewer districts qualify.
Communities that genuinely need help are penalized for having slightly better
neighbors. This is a structural failure of the binary cutoff, not of the
partitioning method.


\subsubsection{Proposed Framework: Continuous Composite Scoring}\label{sec:score-function}

We propose replacing the binary IMU threshold with a continuous underservice
score that serves three roles: ranking districts by need, functioning as the
energy in the FalCom optimization, and allocating funding proportionally rather
than through a binary gate.

\paragraph{Score definition.}
Let $\mathcal{X} = \{x_1, \ldots, x_K\}$ denote a set of $K$ indicators
measured on each census tract $t$. The underservice score is
\begin{equation}\label{eq:score}
  S(t) = \sum_{k=1}^{K} w_k \, \varphi_k\bigl(x_k(t)\bigr),
\end{equation}
where $\varphi_k : \R \to [0,1]$ is a monotone normalization function for
indicator $k$, and $w_k > 0$ are weights satisfying $\sum_k w_k = 1$.

\paragraph{Indicator set.}
The framework accommodates any set of indicators. Table~\ref{tab:indicators}
lists the indicators available in our Chicago data, organized by category.
The original four IMU components are retained for comparability with HRSA,
supplemented by factors the 1970s formula omits.

\begin{table}[H]
\centering
\caption{Candidate indicators for the underservice score.}
\label{tab:indicators}
\begin{tabular}{llll}
\toprule
\textbf{Category} & \textbf{Indicator} & \textbf{Source} & \textbf{Level} \\
\midrule
Access   & CTA transit time to nearest FQHC     & CTA/OTP     & Block  \\
Access   & Provider-to-population ratio          & NPPES+Census & Tract  \\
Coverage & Uninsured rate                        & Chicago Health Atlas & Comm.\ area \\
Burden   & Diabetes prevalence                   & Chicago Health Atlas & Comm.\ area \\
Burden   & Hypertension prevalence               & Chicago Health Atlas & Comm.\ area \\
Demog.   & Poverty rate                          & ACS 5-Year  & Tract  \\
Demog.   & \% aged 65+                           & Decennial   & Block  \\
Demog.   & Infant mortality rate                 & Chicago Health Atlas & Comm.\ area \\
\bottomrule
\end{tabular}
\end{table}

\paragraph{Normalization.}
We use two normalization functions depending on the indicator type.

For indicators with a natural policy threshold (e.g., 30 minutes of transit
time), we use a sigmoid:
\begin{equation}\label{eq:sigmoid}
  \varphi_k(x) = \frac{1}{1 + e^{-\kappa_k (x - \mu_k)}},
\end{equation}
where $\mu_k$ is the threshold (e.g., $\mu = 30$ for transit minutes) and
$\kappa_k$ controls the steepness. This produces a smooth inflection at the
policy-relevant boundary.

For all other indicators we use the percentile rank:
\begin{equation}\label{eq:percentile}
  \varphi_k(x) = \frac{\operatorname{rank}(x)}{N},
\end{equation}
where $N$ is the number of tracts. This is robust to outliers and directly
interpretable: a tract scoring 0.85 is in the 85th percentile of need on that
indicator. The CDC Social Vulnerability Index uses an analogous approach
\cite{CDC_SVI}.

\paragraph{Weighting.}
Equal weights ($w_k = 1/K$) serve as the default. The framework supports
user-adjustable weights, allowing the decision maker to specify which
factors matter and by how much. This separates the policy question (what
should count as underserved?) from the optimization question (given a
definition, where should facilities go?). The model finds the best
configuration for any weight profile.

\paragraph{Dynamic and static components.}
Among the indicators in Table~\ref{tab:indicators}, transit time to the
nearest FQHC is the only component that changes when a facility is added or
removed. All other indicators are fixed characteristics of the tract. This
means the score responds to the decision variable: when FalCom places a
facility in a district, the transit component of $S(t)$ decreases for nearby
tracts, lowering the district energy. The optimizer can observe how each
candidate location improves the system-wide score.


\subsubsection{Probabilistic Extension}\label{sec:prob-imu}

Two of the four original IMU indicators come from the American Community
Survey, where in approximately 50\% of census tracts the margin of error
exceeds 50\% of the estimate \cite{spielman2014patterns}. Rather than
computing a single deterministic score, we model each ACS-sourced indicator
with its reported uncertainty distribution and produce a probability of
underservice for each tract.

For each tract $t$, let $\hat{x}_k(t)$ and $\sigma_k(t)$ denote the ACS
point estimate and standard error for indicator $k$. We draw $M$ samples
\begin{equation}\label{eq:prob-imu}
  x_k^{(m)}(t) \sim \mathcal{N}\bigl(\hat{x}_k(t),\, \sigma_k(t)^2\bigr),
  \qquad m = 1, \ldots, M,
\end{equation}
compute the score $S^{(m)}(t)$ for each sample, and report
\begin{equation}\label{eq:prob-underserved}
  P(\text{underserved} \mid t) =
    \frac{1}{M} \sum_{m=1}^{M}
    \mathbf{1}\bigl[S^{(m)}(t) > \tau\bigr],
\end{equation}
where $\tau$ is any threshold of interest. A tract with a point estimate just
above $\tau$ but high variance may have only a 40\% probability of truly
exceeding it. This replaces the hard cutoff with a continuous risk measure.

For the FalCom energy function, the expected score $\E[S(t)]$ can be used
directly, or a risk-averse variant
$\E[S(t)] + \lambda \sqrt{\Var[S(t)]}$
that penalizes uncertainty, directing resources toward tracts where either
need is confidently high or measurement uncertainty is itself a problem
warranting investment in better data.


\subsubsection{District-Level Aggregation}\label{sec:district-score}

FalCom optimizes over districts, not individual tracts. The energy of a
district $d$ with assigned facility $j$ is the population-weighted average
of its tracts' scores, modulated by the travel time to the facility:
\begin{equation}\label{eq:district-energy}
  E(d) = \frac{1}{\operatorname{Pop}(d)}
    \sum_{t \in d} \operatorname{Pop}(t) \cdot S(t) \cdot
    g\bigl(\tau(t, j)\bigr),
\end{equation}
where $\tau(t, j)$ is the CTA transit time from tract $t$ to facility $j$
and $g$ is a monotone increasing function (e.g., $g(x) = 1 + x/30$) so that
tracts far from their assigned facility contribute more energy. This combines
need ($S$) with access ($g$) in a single objective that FalCom minimizes
across the system.


\subsubsection{Policy Design: Continuous Funding Allocation}\label{sec:transition}

Any improved scoring methodology faces the political constraint described in
Section~\ref{sec:imu-political}: existing FQHCs will oppose any formula change
that threatens their funding. We propose a transition mechanism that evaluates
all sites under the new score without defunding any existing clinic.

\paragraph{Freeze the floor, redirect the growth.}
In the transition year, every existing FQHC receives a funding floor equal to
its current grant amount. The continuous score is computed for all districts.
In subsequent years, the floor remains nominally frozen: no site receives less
than it received the previous year. All new appropriations above the aggregate
floor are allocated proportionally to the district scores. Sites in the
highest-need districts receive the growth; sites in lower-need districts
receive flat nominal funding.

Federal FQHC appropriations have grown from \$554 million in 1990 to over
\$6 billion today. Inflation alone erodes the frozen floor by 2--3\% per year
in real terms. Over a 10--15 year horizon, a site whose district has genuinely
improved receives the same nominal amount but substantially less in real
dollars, while high-need districts grow. No clinic is defunded. No
designation is revoked. The growth is redirected.

\paragraph{Annual review cycle.}
Each year, the partition is recomputed with updated data (new providers,
population shifts, new or closed sites). The score provides HRSA with two
outputs:
\begin{enumerate}
  \item Where the next new facility should go (the district with the highest
    energy after optimization).
  \item Which existing districts have improved enough that their growth
    funding can be reallocated (gradual rebalancing).
\end{enumerate}
This makes the scoring methodology reformable. Because no single annual update
can cause a funding cliff, the formula can be improved incrementally as better
data or indicators become available.


\subsection{Referral System Design}

\textcolor{blue}{TODO: Each FQHC must be within a certain CTA travel time of a referral hospital. Referral constraints. Integration with the hierarchical model.}



%%%%%%%%%%%%%%%%%%%%%%%%%%%%%%%%%%%%%%%%%%%%%%%%%%
\section{Experiments}\label{sec:experiments}
%%%%%%%%%%%%%%%%%%%%%%%%%%%%%%%%%%%%%%%%%%%%%%%%%%

\subsection{Data Sources}

\paragraph{Census and demographic data.}
The newest block-level data is from the 2020 Decennial Census (DHC). Block-level demographic data is only released with the 
decennial census (every 10 years). Poverty rate data comes from ACS 5-Year 2020 (mean 19.0\% at tract level). Total Chicago 
population: 2,783,891; population 65+: 356,716.

\paragraph{Provider data.}
We use NPPES Data Dissemination V.2 (February 2026), which contains 77,590 Chicago providers from 9.3 million total records nationwide. 
This includes all provider types; we filter for primary care physicians.

\paragraph{FQHC demand and capacity (UDS).}
The Uniform Data System (UDS) collects annual reports from every HRSA-funded health center containing detailed information about 
patients served, medical/dental/behavioral visits, demographics, workforce (physicians, nurse practitioners), services, and finance. 
This dataset provides facility capacity, actual patient demand, service utilization, and staffing constraints. Almost no operations 
research papers use this data.

\paragraph{Area Health Resource File (AHRF).}
This dataset covers every U.S.\ county and includes hundreds of variables: physician counts, hospital beds, demographics, poverty rates, 
infant mortality, insurance coverage, and workforce data.

\begin{table}[h]
\centering
\caption{Chicago geographic and provider data files.}
\label{tab:chicago-data}
\scalebox{0.8}{
\begin{tabular}{llll}
\toprule
\textbf{File} & \textbf{Level} & \textbf{Year} & \textbf{Contents} \\
\midrule
chicago\_blocks\_enriched.pkl & Block & 2020 & Population, age 65+, poverty rate \\
chicago\_block\_age65.csv & Block & 2020 DHC & GEOID20, total population, population 65+ \\
chicago\_tract\_poverty.csv & Tract & 2020 ACS & Poverty rate, population below poverty \\
chicago\_providers.csv & Provider & 2026 NPPES & 77,590 Chicago NPI providers \\
\bottomrule
\end{tabular}
}
\end{table}


\subsection{Implementation}

\textcolor{blue}{TODO: FalCom configuration for Chicago. CTA travel time computation. Dashboard description.}


\subsection{Results}

\textcolor{blue}{TODO}



%%%%%%%%%%%%%%%%%%%%%%%%%%%%%%%%%%%%%%%%%%%%%%%%%%
\section{Discussion}\label{sec:discussion}
%%%%%%%%%%%%%%%%%%%%%%%%%%%%%%%%%%%%%%%%%%%%%%%%%%

\textcolor{blue}{TODO}


%%%%%%%%%%%%%%%%%%%%%%%%%%%%%%%%%%%%%%%%%%%%%%%%%%
\section{Conclusion}\label{sec:conclusion}
%%%%%%%%%%%%%%%%%%%%%%%%%%%%%%%%%%%%%%%%%%%%%%%%%%

\textcolor{blue}{TODO}


\bibliographystyle{plainnat}
\bibliography{ref}

\end{document}
